% Answers to Chapter 16, translated from sv/back/facit/16.tex.

\facitavsnitt{c16}{Chapter 16}{Complex numbers, part II}
\begin{facit}

\svar{1.1} $I \approx 5\,\angle\,0.93\,\text{mA}$, with the angle $0.93\,\text{rad} \approx
53.1^{\circ}$.

\svar{1.2} $U = \sqrt{2} + j\sqrt{2} \approx 1.41 + j1.41\,\text{V}$

\svar{1.3} $Z \approx 0.4\,\angle\,{-0.14}\,\text{k}\Omega$ (the angle in radians), in
rectangular form $Z \approx 0.4 - j0.06\,\text{k}\Omega$; more precisely $396 - j57\,\Omega$.

\svar{2.1} \sv{a} $a = 4 + j3$, $b = -3 + j6$, $c = -1 + j10$
\sv{b} Arrows from the origin to $(4, 3)$, $(-3, 6)$ and $(-1, 10)$: $a$ in the first quadrant,
$b$ and $c$ in the second.
\sv{c} $\abs{a} = 5$, $\abs{b} = \sqrt{45} \approx 6.7$, $\abs{c} = \sqrt{101} \approx 10.0$
\sv{d} $b + 2c = -5 + j26$, $\abs{b + 2c} = \sqrt{701} \approx 26.5$
\sv{e} $\delta_a \approx 0.64\,\text{rad} \approx 36.9^{\circ}$,
$\delta_b \approx 2.0\,\text{rad} \approx 116.6^{\circ}$,
$\delta_c \approx 1.7\,\text{rad} \approx 95.7^{\circ}$
\sv{f} $d = -8 - j6$

\svar{2.2} \sv{a} An arrow from the origin to $(-12, 4)$, in the second quadrant.
\sv{b} $\abs{U} = \sqrt{160} \approx 12.6\,\text{V}$, $\delta \approx 2.8\,\text{rad} \approx
161.6^{\circ}$, so $U \approx 12.6\,e^{j2.8}\,\text{V}$.
\sv{c} $\omega = 200\pi \approx 628\,\text{rad/s}$
\sv{d} $u(t) \approx 12.6\,e^{j(200\pi t + 2.8)}\,\text{V}$

\svar{3.1} \sv{a} $j$ \sv{b} $-1$ \sv{c} $1$ \sv{d} $-j$ \sv{e} $1$

\svar{3.2} \sv{a} $4$ \sv{b} $2.5 + j4.33$ \sv{c} $\approx 1.41 - j1.41$
\sv{d} $-5 + j8.66$

\svar{3.3} \sv{a} $3\sqrt{2}\,e^{j\pi/4} \approx 4.24\,e^{j\pi/4}$
\sv{b} $2\sqrt{2}\,e^{j3\pi/4} \approx 2.83\,e^{j3\pi/4}$
\sv{c} $6\,e^{-j\pi/2}$ \sv{d} $5\,e^{j\pi}$

\svar{3.4} \sv{a} $12\,\angle\,75^{\circ}$
\sv{b} $\frac{4}{3}\,\angle\,{-15^{\circ}} \approx 1.33\,\angle\,{-15^{\circ}}$
\sv{c} $16\,\angle\,60^{\circ}$
\sv{d} The magnitudes are multiplied and the angles added; rectangular form needs four partial
products and the handling of $j^2$.

\svar{3.5} \sv{a} $-1 + j7$
\sv{b} $z_1 \approx 5\,\angle\,53.1^{\circ}$, $z_2 = \sqrt{2}\,\angle\,45^{\circ} \approx
1.41\,\angle\,45^{\circ}$
\sv{c} $\approx 7.07\,\angle\,98.1^{\circ}$
\sv{d} Yes: $\abs{-1 + j7} = \sqrt{50} \approx 7.07$ and the angle is $\approx 98.1^{\circ}$.

\svar{3.6} \sv{a} $5\,\angle\,60^{\circ}$ \sv{b} $3\,\angle\,{-90^{\circ}}$
\sv{c} $2\,\angle\,{-90^{\circ}}$
\sv{d} It decreases by $90^{\circ}$ and the magnitude is unchanged, since
$j = 1\,\angle\,90^{\circ}$: the number is rotated a quarter turn clockwise.

\svar{3.7} \sv{a} $4\,\angle\,60^{\circ}$ \sv{b} $8\,\angle\,90^{\circ} = j8$
\sv{c} $16\,\angle\,120^{\circ}$ \sv{d} $z^n = \abs{z}^n\,\angle\,n\delta$ (de Moivre's formula)

\svar{3.8} \sv{a} $10\,\angle\,\frac{\pi}{6} = 10\,\angle\,30^{\circ}$
\sv{b} $4\,\angle\,{-\frac{\pi}{2}} = 4\,\angle\,{-90^{\circ}}$
\sv{c} $6\,\angle\,0^{\circ}$
\sv{d} The signals must have the same angular frequency $\omega$.

\svar{3.9} \sv{a} $u(t) = 8\sin\left(100\pi t + \frac{\pi}{4}\right)\,\text{V}$
\sv{b} $u(t) = 5\sin\left(100\pi t - \frac{\pi}{2}\right)\,\text{V}$
\sv{c} $i(t) \approx 5\sin(100\pi t + 0.927)\,\text{mA}$
\sv{d} $i(t) \approx 2.83\sin(100\pi t - 2.36)\,\text{mA}$

\svar{3.10} \sv{a} $U_1 = 6 + j0\,\text{V}$, $U_2 = 0 + j8\,\text{V}$ \sv{b} $6 + j8\,\text{V}$
\sv{c} $10\,\angle\,53.1^{\circ}\,\text{V}$
\sv{d} $u(t) \approx 10\sin(100\pi t + 0.927)\,\text{V}$

\svar{3.11} \sv{a} $Z = 5\,\angle\,45^{\circ}\,\Omega$ \sv{b} $Z \approx 3.54 + j3.54\,\Omega$
\sv{c} Inductive: the imaginary part is positive.
\sv{d} $R \approx 3.54\,\Omega$ and $X \approx 3.54\,\Omega$

\svar{3.12} \sv{a} $Z_L = j20\,\Omega$ \sv{b} $Z_C = -j20\,\Omega$ \sv{c} $Z = 15\,\Omega$
\sv{d} Resonance: the reactances cancel, $\omega L = 1/(\omega C)$, and the impedance is purely
resistive.

\svar{3.13} \sv{a} $a = 5 + j12$, $b = -8 + j6$ \sv{b} $\abs{a} = 13$, $\abs{b} = 10$
\sv{c} $a + b = -3 + j18$, $\abs{a + b} = \sqrt{333} \approx 18.25$
\sv{d} $\delta_a \approx 67.4^{\circ}$, $\delta_b \approx 143.1^{\circ}$

\svar{3.14} \sv{a} $jz = -4 + j3$ \sv{b} Equal: $\abs{z} = \abs{jz} = 5$.
\sv{c} $53.1^{\circ}$ and $143.1^{\circ}$; the difference is $90^{\circ}$.
\sv{d} It rotates the number $90^{\circ}$ anticlockwise without changing its length.

\svar{3.15} \sv{a} True: $\abs{e^{j\delta}} = \sqrt{\cos^2\delta + \sin^2\delta} = 1$.
\sv{b} False: the magnitudes are multiplied, $\abs{z_1z_2} = \abs{z_1}\abs{z_2}$.
\sv{c} True: $z_1/z_2 = (\abs{z_1}/\abs{z_2})\,\angle\,(\delta_1 - \delta_2)$.
\sv{d} False: a phasor carries only amplitude and phase; $\omega$ is given separately.
\sv{e} True: it is Euler's identity, $e^{j\pi} = -1$.
\sv{f} False: multiplication is easiest in polar form or Euler form.

\svar[Test yourself]{T} \sv{1} $Z = 5\sqrt{3} + j5 \approx 8.66 + j5$
\sv{2} $e^{j\pi} = -1$: the number $1$ is rotated half a turn ($180^{\circ}$) along the unit
circle and ends up at $-1$.

\end{facit}
